\chapter{Herbal Medicine}

\section*{Definition and Overview}
Herbal medicine is the use of plant-derived products, such as leaves, roots, seeds, and extracts, for therapeutic purposes. It is one of the oldest forms of medicine and is still widely used, both as traditional practice and as commercial products such as tablets, teas, and creams. While some herbal remedies have genuine, well-documented benefits, many have limited or no evidence of effectiveness, and some interact dangerously with conventional medicines. Using herbal products safely means informing health professionals about them, buying quality products, and understanding that "natural" does not mean risk-free.

\section*{Epidemiology}
Complementary medicines, including herbal products, are used by around half of Australians at some point, and Australians spend more than \$3 billion a year on them. Herbal products are most commonly used for colds, pain, sleep, anxiety, digestive complaints, and menopausal symptoms. Many people use them alongside prescription medicines, and a substantial proportion do not tell their doctor. While most herbal products cause no harm, adverse events and dangerous drug interactions occur, and this under-reporting is a recognised safety problem.

\section*{Aetiology and Pathophysiology}
\begin{itemize}
    \item \textbf{Active constituents:} plants contain biologically active chemicals, some of which are the source of modern drugs, such as digoxin from foxglove and aspirin from willow bark
    \item \textbf{Mechanisms:} herbal products act through a wide range of pharmacological effects, depending on the plant and preparation
    \item \textbf{Standardisation:} quality varies enormously; products differ in potency, purity, and contaminants
    \item \textbf{Regulation:} in Australia, complementary medicines are regulated by the TGA, but lower-risk products are not assessed for efficacy
    \item \textbf{Well-supported uses:} a few herbs, such as St John's wort for mild depression, have evidence of benefit
    \item \textbf{Limited evidence:} most herbal claims are supported by weak or conflicting studies
    \item \textbf{Risks:} contamination, mislabelling, liver and kidney toxicity, allergic reactions, and interactions with medicines
\end{itemize}

\section*{Clinical Features}
\begin{itemize}
    \item People commonly seek herbal remedies for chronic and difficult-to-treat symptoms
    \item Use may be reported as tea, tablets, tinctures, creams, or traditional preparations
    \item Improvement is often subjective and influenced by expectation
    \item Adverse effects may be mild, such as digestive upset, or serious, such as liver injury
    \item Interactions can change the effect of warfarin, antidepressants, immunosuppressants, and other drugs
    \item Herbal products can cause unexplained symptoms that are only traced to the product with a careful history
\end{itemize}

\section*{Differential Diagnosis}
\begin{itemize}
    \item Symptoms attributed to herbal remedies may actually reflect the underlying illness
    \item Liver or kidney injury from herbal products must be distinguished from other causes
    \item Worsening of a condition may result from abandoning effective treatment in favour of herbs
    \item Interactions with prescribed medicines can cause new symptoms
    \item A thorough medicine and complementary medicine history is essential in unexplained illness
\end{itemize}

\section*{When to Seek Medical Care}
Always tell your doctor and pharmacist about any herbal products you use, especially before surgery, during pregnancy or breastfeeding, or if you take other medicines. Seek medical review for any new symptom after starting a herbal product.

\textbf{Call 000 (or local emergency services) for:}
\begin{itemize}
    \item Jaundice, dark urine, or severe abdominal pain (possible liver injury)
    \item Chest pain, palpitations, or collapse
    \item Severe allergic reaction: swelling, difficulty breathing, or rash
    \item Severe bleeding or easy bruising in someone taking blood thinners and herbal products
    \item Fits, confusion, or loss of consciousness
\end{itemize}

\section*{Diagnosis and Investigations}
\begin{itemize}
    \item \textbf{Full medicine history:} includes all herbal, vitamin, and complementary products
    \item \textbf{Liver and kidney tests:} when toxicity is suspected
    \item \textbf{Product identification:} the exact product and dose should be documented, including brand and batch
    \item \textbf{Interaction review:} pharmacists can check for interactions with prescribed medicines
    \item \textbf{Reporting:} adverse reactions to complementary medicines should be reported to the TGA
    \item \textbf{Investigations:} guided by the symptoms and suspected cause
\end{itemize}

\section*{Management}
\begin{itemize}
    \item \textbf{Evidence-based discussion:} health professionals review what is known about the product's benefits, risks, and interactions
    \item \textbf{Quality products:} if herbs are used, choose reputable brands with TGA registration and standardised extracts
    \item \textbf{Avoid high-risk situations:} pregnancy, breastfeeding, children, surgery, and people with liver or kidney disease
    \item \textbf{Stop if problems:} cease the product and seek advice if any adverse effect occurs
    \item \textbf{Interaction management:} adjust prescriptions under medical supervision where interactions exist
    \item \textbf{Integration:} herbal medicine is best used in discussion with the treating team rather than as a substitute for proven care
\end{itemize}

\section*{Complications}
\begin{itemize}
    \item Liver injury from products such as some traditional Chinese medicines and kava
    \item Kidney damage and electrolyte disturbance
    \item Dangerous interactions with anticoagulants, antidepressants, and other drugs
    \item Interference with anaesthesia and surgery
    \item Contamination with heavy metals, pesticides, or undeclared drug ingredients
    \item Delayed or abandoned effective treatment
    \item Allergic and toxic reactions
\end{itemize}

\section*{Prognosis and Outcomes}
The outcomes of herbal medicine use depend entirely on the product and the condition. A small number of herbs have evidence of benefit for specific indications, while most have uncertain or no proven value. When used with medical knowledge and for well-supported purposes, harm is uncommon; when used in place of proven treatment or in ignorance of interactions, outcomes can be poor. An open, informed partnership between patients and health professionals gives the best results.

\section*{Prevention and Self-Care}
\begin{itemize}
    \item Tell your doctor and pharmacist about all herbal products
    \item Ask about evidence for effectiveness before buying a product
    \item Choose TGA-registered products from reputable suppliers
    \item Avoid herbal products in pregnancy, breastfeeding, and young children unless advised by a qualified practitioner
    \item Stop herbal products before planned surgery after discussing with your team
    \item Beware of products promising cures for serious disease
    \item Report any adverse reaction to the TGA
    \item Be cautious combining herbs with blood thinners, antidepressants, and immunosuppressants
\end{itemize}

\section*{Sources and Further Reading}
The following reputable sources provide further information for healthcare students and patients seeking reliable, current advice about herbal medicine.

\begin{itemize}
    \item Therapeutic Goods Administration. \textit{Complementary medicines regulation and safety.} https://www.tga.gov.au/
    \item Healthdirect. \textit{Complementary medicines and herbal products.} https://www.healthdirect.gov.au/complementary-medicines
    \item National Centre for Complementary and Integrative Health. \textit{Herbal medicines: what you should know.} https://www.nccih.nih.gov/
    \item Better Health Channel. \textit{Herbal medicine.} https://www.betterhealth.vic.gov.au/
    \item MedlinePlus. \textit{Herbal medicine.} https://medlineplus.gov/herbalmedicine.html
    \item NHS. \textit{Herbal medicines: safety and regulation.} https://www.nhs.uk/conditions/complementary-therapies/
\end{itemize}
